\chapter{核心代码附录}

\section{ResNet-50模型训练代码}

以下为基于MindSpore框架的ResNet-50模型训练核心代码：

\begin{codeblock}[language=Python, caption={ResNet-50模型训练代码}]
import mindspore as ms
from mindspore import nn, context, Tensor
from mindspore.train import Model, LossMonitor
from mindspore.dataset import vision, transforms
from mindspore.dataset import Cifar10Dataset

context.set_context(mode=context.GRAPH_MODE, device_target="GPU")

# 数据预处理
transform_train = transforms.Compose([
    vision.RandomCrop(32, padding=4),
    vision.RandomHorizontalFlip(),
    vision.Resize(224),
    vision.Normalize(mean=[0.4914, 0.4822, 0.4465],
                     std=[0.2023, 0.1994, 0.2010]),
    vision.HWC2CHW()
])

# 数据集加载
dataset = Cifar10Dataset(dataset_dir="./cifar-10",
    usage="train", shuffle=True)
dataset = dataset.map(transform_train,
    input_columns="image")
dataset = dataset.batch(128)

# 模型定义
net = resnet50(num_classes=10)
loss_fn = nn.CrossEntropyLoss()
opt = nn.SGD(net.trainable_params(),
    learning_rate=0.1, momentum=0.9,
    weight_decay=5e-4)
model = Model(net, loss_fn, opt)

# 模型训练
model.train(200, dataset,
    callbacks=[LossMonitor(0.1)],
    dataset_sink_mode=True)
\end{codeblock}

\section{INT8量化转换代码}

以下为模型INT8后训练量化的核心代码：

\begin{codeblock}[language=Python, caption={INT8量化转换代码}]
from mindspore_lite import Converter

converter = Converter()
converter.set_config_file("./quantization.cfg")
converter.set_model_file("resnet50.mindir")
converter.set_output_file("resnet50_quant.ms")
converter.set_weight_quant(True)
converter.set_input_shape("input:1,3,224,224")
converter.convert()
\end{codeblock}

\section{鸿蒙NAPI推理接口代码}

以下为鸿蒙端侧推理引擎的NAPI接口核心代码：

\begin{codeblock}[language=C++, caption={NAPI推理接口代码}]
#include <mindspore/lite/model.h>
#include <mindspore/lite/context.h>

napi_value Predict(napi_env env, napi_callback_info info) {
    // 获取输入数据
    size_t argc = 2;
    napi_value args[2];
    napi_get_cb_info(env, info, &argc, args, nullptr, nullptr);

    // 创建推理上下文
    auto context = std::make_shared<mindspore::lite::Context>();
    context->device_list_ = {{mindspore::lite::DT_CPU, {}}};
    context->cpu_context_ = {
        .enable_float16_ = false,
        .thread_num_ = 4
    };

    // 创建模型并加载
    auto model = mindspore::lite::Model::CreateModel(
        model_path, context);

    // 创建推理器
    auto runner = mindspore::lite::ModelParallelRunner();
    runner.Init(model_path, context);

    // 执行推理
    auto inputs = runner.GetInputs();
    CopyInputData(inputs, input_buffer);
    auto outputs = runner.Predict(inputs);

    // 返回分类结果
    napi_value result;
    napi_create_string_utf8(env,
        GetTopResult(outputs).c_str(),
        NAPI_AUTO_LENGTH, &result);
    return result;
}
\end{codeblock}

\section{ArkTS界面核心代码}

以下为鸿蒙ArkTS应用主界面的核心代码：

\begin{codeblock}[caption={ArkTS主界面代码}]
@Entry
@Component
struct IndexPage {
  @State imagePath: string = ''
  @State resultText: string = ''
  @State confidenceList: number[] = []

  build() {
    Column() {
      // 图片展示区域
      Image(this.imagePath)
        .width(300).height(300)
        .objectFit(ImageFit.Cover)
        .borderRadius(8)

      // 分类按钮
      Button('开始分类')
        .width('80%').height(40)
        .onClick(() => this.classifyImage())

      // 分类结果展示
      Text(this.resultText)
        .fontSize(20).fontWeight(FontWeight.Bold)

      // 置信度柱状图
      ForEach(this.confidenceList,
        (item: number, index: number) => {
          Row() {
            Progress({ value: item, total: 100 })
              .width('60%').color('#007DFF')
          }
        })
    }.width('100%').padding(20)
  }

  async classifyImage() {
    let result = await mindsporeLite.predict(
      this.imagePath)
    this.resultText = result.topClass
    this.confidenceList = result.confidences
  }
}
\end{codeblock}